\section*{Chapitre \ref{ch:b2:aromatic-substitution} --- Aromaticité et substitution électrophile aromatique}

\begin{solution}{exo:b2:aromatic-substitution:1}
\ce{HNO3 + 2H2SO4 -> NO2+ + H3O+ + 2HSO4-}~: l'acide nitrique est protoné, puis perd de l'eau.
\end{solution}

\begin{solution}{exo:b2:aromatic-substitution:2}
Le toluène donne le 2-bromotoluène et le 4-bromotoluène~: le méthyle oriente vers les positions
ortho et para. Le nitrobenzène donne le 3-bromonitrobenzène, plus lentement~: le groupe nitro
oriente en méta.
\end{solution}

\begin{solution}{exo:b2:aromatic-substitution:3}
\ce{-CH3}~: activant, o/p. \ce{-OH}~: fortement activant, o/p. \ce{-Cl}~: désactivant, o/p.
\ce{-NO2}~: fortement désactivant, méta. \ce{-COCH3}~: désactivant, méta. \ce{-NHCOCH3}~:
activant (modérément), o/p.
\end{solution}

\begin{solution}{exo:b2:aromatic-substitution:4}
L'acétophénone (1-phényléthan-1-one), \ce{C6H5COCH3}, par l'intermédiaire de l'\omterm{def:b2:aromatic-substitution:friedel-crafts}{ion acylium} \ce{CH3CO+}.
\end{solution}

\begin{solution}{exo:b2:aromatic-substitution:5}
Le carbocation primaire (ou son complexe) se réarrange par migration d'hydrure en cation
isopropyle secondaire, qui alkyle le cycle. Propylbenzène~: \omterm{def:b2:aromatic-substitution:friedel-crafts}{acylation de Friedel--Crafts} par le
chlorure de propanoyle (sans réarrangement), puis réduction du \ce{C=O} de la cétone en \ce{CH2}.
\end{solution}

\begin{solution}{exo:b2:aromatic-substitution:6}
Le groupe \ce{-OH} est un donneur mésomère fort~: le cycle est si activé que le dibrome
suffit comme électrophile, et les trois positions ortho et para sont toutes substituées.
\end{solution}

\begin{solution}{exo:b2:aromatic-substitution:7}
3-Bromonitrobenzène~: nitrer, puis bromer (le nitro oriente en méta). 4-Bromonitrobenzène~:
bromer, puis nitrer (le brome oriente en ortho/para), et séparer l'isomère para de l'isomère
ortho.
\end{solution}

\begin{solution}{exo:b2:aromatic-substitution:8}
Le mélange nitrant oxyde l'aniline, très riche en électrons, et protone son azote~:
\ce{-NH3+} est un \omterm{def:b2:aromatic-substitution:directing}{groupe désactivant}, \omterm{def:b2:aromatic-substitution:directing}{orienteur méta}. Dans l'acétanilide, le doublet de l'azote est
partagé avec le carbonyle~: moins basique et moins activant, l'azote n'est ni protoné ni oxydé,
et oriente toujours en ortho/para, surtout en para pour des raisons stériques.
\end{solution}

\begin{solution}{exo:b2:aromatic-substitution:9}
Sulfoner le toluène (surtout en para, le groupe volumineux évitant l'ortho)~; bromer~: le brome entre
en ortho du méthyle (la position para est bloquée, et les deux groupes favorisent cette position)~;
retirer le groupe \ce{SO3H} par chauffage en milieu acide dilué.
\end{solution}

\begin{solution}{exo:b2:aromatic-substitution:10}
Les deux groupes orientent en ortho/para~; leurs positions para sont occupées l'une par l'autre. Le méthoxy,
l'activant le plus fort, l'emporte~: nitration en ortho de \ce{-OCH3} (position 2), qui donne le
4-méthyl-2-nitroanisole.
\end{solution}

\begin{solution}{exo:b2:aromatic-substitution:11}
Chaque groupe nitro désactive fortement le cycle~: la deuxième nitration exige des conditions plus
dures que la première, et la troisième (sur un cycle portant deux groupes nitro et seulement le méthyle,
faiblement activant) des conditions plus dures encore.
\end{solution}

\begin{solution}{exo:b2:aromatic-substitution:12}
Acide 4-nitrobenzoïque~: nitrer le toluène, séparer l'isomère para, oxyder le méthyle en
\ce{-COOH} (permanganate à chaud). Acide 3-nitrobenzoïque~: oxyder d'abord le toluène en acide benzoïque, puis
nitrer (\ce{-COOH} oriente en méta).
\end{solution}

\begin{solution}{pb:b2:aromatic-substitution:1}
\textbf{1.} \ce{HNO3 + 2H2SO4 -> NO2+ + H3O+ + 2HSO4-}.
\textbf{2.} \ce{O=N+=O}~: linéaire, isoélectronique de \ce{CO2}.
\textbf{3.} L'azote se lie au carbone du cycle situé en para de \ce{CH3}, qui devient tétraédrique
(et porte H et \ce{NO2}).
\textbf{4.} La charge positive se trouve sur les deux carbones en ortho du carbone attaqué et sur le
carbone en para de celui-ci, qui porte le méthyle.
\textbf{5.} La première, la formation de l'\omterm{def:b2:aromatic-substitution:seAr}{intermédiaire de Wheland}.
\textbf{6.} Une base du milieu~: \ce{HSO4-} ou l'eau.
\textbf{7.} Deux ortho, deux méta, une para.
\textbf{8.} 2 : 2 : 1, soit 40~\% ortho, 40~\% méta, 20~\% para.
\textbf{9.} Pour une attaque en méta, la charge positive ne se trouve jamais sur le carbone qui porte le méthyle~;
aucune forme n'est aidée par le méthyle.
\textbf{10.} Une forme mésomère porte la charge sur le carbone lié à \ce{CH3}~: un carbocation
tertiaire, stabilisé par l'effet inductif et l'\omterm{def:b2:fragment-orbitals:hyperconjugation}{hyperconjugaison} du méthyle.
\textbf{11.} L'ortho et le para sont favorisés (96~\% contre 60~\% au hasard)~; le méta est presque absent.
\textbf{12.} Par position~: ortho $58/2 = 29~\%$, para 38~\%~: le para est la position la plus réactive,
les positions ortho étant légèrement encombrées par le méthyle.
\textbf{13.} Plus vite~: le méthyle active le cycle.
\textbf{14.} Le 4-nitrotoluène (fusion à \qty{52}{\degreeCelsius}).
\textbf{15.} Refroidir le mélange brut~: l'isomère para cristallise et est séparé par filtration~; le
liquide retient les isomères ortho et méta.
\textbf{16.} Une molécule symétrique s'empile mieux dans un cristal~: plus d'interactions par unité de volume,
donc une température de fusion plus élevée.
\textbf{17.} Par \omterm{def:b2:liquid-vapour-diagrams:fractional-distillation}{distillation fractionnée} sous pression réduite, ou par chromatographie~; ou par
un refroidissement plus poussé (l'isomère méta fond à \qty{16}{\degreeCelsius}).
\textbf{18.} Les isomères ont des températures d'ébullition voisines (même formule, polarité semblable).
\textbf{19.} \ce{C7H8}~: \qty{92.14}{g/mol}~; \ce{C7H7NO2}~: \qty{137.14}{g/mol}.
\textbf{20.} $100/92.14 = \qty{1.085}{mol}$ de toluène~; $0.90 \times 1.085 = \qty{0.977}{mol}$ de
nitrotoluènes.
\textbf{21.} Total $0.977 \times 137.14 = \qty{134.0}{g}$~: ortho \qty{77.7}{g}, méta
\qty{5.4}{g}, para \qty{50.9}{g}.
\textbf{22.} Une seconde nitration (dinitrotoluènes), et l'oxydation du groupe méthyle.
\textbf{23.} Oxyder le méthyle en \ce{-COOH} par le permanganate en milieu basique à chaud, puis acidifier.
\textbf{24.} L'opération donne $\boldsymbol{\approx \qty{51}{g}}$ de 4-nitrotoluène.
\end{solution}
